
\setcounter{chapter}{18}

\chapter{特征标理论的例子与应用}

\section{小阶群的特征标}

有限群 \(G\) 的\textbf{特征标表}是按如下方式排版的表：把 \(r\) 个共轭类的代表元列在顶行，把不可约特征标列在第一列。表中位于行 \(\chi_i\) 与列 \(g_j\) 的项是 \(\chi_i(g_j)\).

有限群的特征标表在行的置换与列的置换下是唯一的。习惯上把主特征标放在第一行、把单位元放在第一列，并按次数递增排列特征标。在我们的例子中，我们将在每个类下面列出该共轭类的大小，于是整个表有 \(r+2\) 行与 \(r+1\) 列（尽管严格说来，特征标表是特征标值的 \(r \times r\) 矩阵）。这将使人们容易用第一正交关系检验「行的正交性」：若类由 \(g_1, \ldots, g_r\) 代表、大小分别为 \(d_1, \ldots, d_r\)，则
\[
(\chi_i, \chi_j) = \frac{1}{|G|}\sum_{k=1}^{r}d_k\chi_i(g_k)\overline{\chi_j(g_k)}.
\]
第二正交关系说特征标表的任意两个不同列的 Hermite 积为零（即它给出「列的正交性」）。

Conway、Curtis、Norton、Parker 与 Wilson 的 \textit{Atlas of Finite Groups}（Clarendon Press, 1985）中给出了许多特征标表。其中包括 Monster 单群 \(M\) 的特征标表。群 \(M\) 有 194 个不可约特征标。\(M\) 的非主不可约特征标的最小次数是 196883，最大次数约为 \(2 \times 10^{26}\). 尽管如此，仍可以计算所有这些特征标在 \(M\) 的所有共轭类上的值。

第一个特征标表的例子：设 \(G = \langle x\rangle\) 是 2 阶循环群。则 \(G\) 有 2 个共轭类与两个不可约特征标：
\[
\begin{array}{c|cc}
\text{类} & 1 & x\\ \hline
\text{大小} & 1 & 1\\
\chi_1 & 1 & 1\\
\chi_2 & 1 & -1
\end{array}
\]
（\(Z_2\) 的特征标表。）

这个阿贝尔群的特征标与表示是同一回事，而任何阿贝尔群的不可约表示在第 18.2 节末的例 1 中描述。

类似地，若 \(G = \langle x\rangle\) 是 3 阶循环群，\(\zeta\) 是固定的本原 3 次单位根（故 \(\zeta^2 = \bar\zeta\)），则 \(G\) 的特征标表为：
\[
\begin{array}{c|ccc}
\text{类} & 1 & x & x^2\\ \hline
\text{大小} & 1 & 1 & 1\\
\chi_1 & 1 & 1 & 1\\
\chi_2 & 1 & \zeta & \zeta^2\\
\chi_3 & 1 & \zeta^2 & \zeta
\end{array}
\]
（\(Z_3\) 的特征标表。）

下面构造 \(S_3\) 的特征标表。回忆第 18.2 节例 2 中 \(S_3\) 有 3 个不可约特征标，其特征标值在该例与第 18.3 节末的例 1 中描述。
\[
\begin{array}{c|ccc}
\text{类} & 1 & (1\ 2) & (1\ 2\ 3)\\ \hline
\text{大小} & 1 & 3 & 2\\
\chi_1 & 1 & 1 & 1\\
\chi_2 & 1 & -1 & 1\\
\chi_3 & 2 & 0 & -1
\end{array}
\]
（\(S_3\) 的特征标表。）

下面考虑 \(D_8\)，采用第 18.3 节例 3 的记号。由推论 11，\(D_8\) 有四个 1 次特征标。此外，在例 3 中我们构造了一个不可约的 2 次表示。由于这些表示次数的平方和为 8，这就给出了全部不可约表示（或者，由于有 5 个共轭类，就有 5 个不可约表示）。若用横线表示到换位子商群（即 Klein 四元群）的过渡，则 \(\bar r^2 = \bar s^2 = 1\). 1 次表示（即其特征标）由把生成元 \(s\) 与 \(r\) 映到 \(\pm1\) 算出（而乘积类映到值的乘积）。2 次不可约表示的矩阵已在第 18.3 节例 3 中算出，这个表示的特征标可直接从这些矩阵读出。故 \(D_8\) 的特征标表为：
\[
\begin{array}{c|ccccc}
\text{类} & 1 & r^2 & s & r & sr\\ \hline
\text{大小} & 1 & 1 & 2 & 2 & 2\\
\chi_1 & 1 & 1 & 1 & 1 & 1\\
\chi_2 & 1 & 1 & -1 & 1 & -1\\
\chi_3 & 1 & 1 & 1 & -1 & -1\\
\chi_4 & 1 & 1 & -1 & -1 & 1\\
\chi_5 & 2 & -2 & 0 & 0 & 0
\end{array}
\]
（\(D_8\) 的特征标表。）

下面计算 8 阶四元数群的特征标表。我们采用通常的表现
\[
Q_8 = \langle i, j \mid i^4 = 1,\ j^2 = i^2,\ i^{-1}ji = j^{-1}\rangle
\]
并令 \(k = ij\) 而 \(-1 = i^2\). \(Q_8\) 的共轭类由 \(1, -1, i, j, k\) 代表，大小分别为 1, 1, 2, 2, 2. 由于 \(Q_8\) 的换位子商群是 Klein 四元群，有四个 1 次特征标。剩下的那一个不可约特征标的次数必为 2，以使次数平方和为 8. 设 \(\chi_5\) 是 \(Q_8\) 的 2 次不可约特征标。人们可以验证第 18.1 节第二组例子中例 7 显式描述的表示 \(\varphi : Q_8 \to GL_2(\mathbb{C})\) 提供 \(\chi_5\)，但我们要说明正交关系如何在并不知道这些显式矩阵的情况下给出 \(\chi_5\) 的值。若 \(\varphi\) 是不可约的 2 次表示，则由 Schur 引理（参看第 18.1 节习题 18），\(\varphi(-1)\) 是 \(2 \times 2\) 纯量矩阵；又因为 \(-1\) 在 \(Q_8\) 中的阶为 2，它是 \(\pm\) 单位矩阵。故 \(\chi_5(-1) = \pm2\).

设 \(\chi_5(i) = a\)、\(\chi_5(j) = b\) 与 \(\chi_5(k) = c\). 正交关系给出
\[
1 = (\chi_5, \chi_5) = \frac{1}{8}\left(2^2+(\pm2)^2+2a\bar a+2b\bar b+2c\bar c\right).
\]
由于 \(a\bar a\)、\(b\bar b\) 与 \(c\bar c\) 是非负实数，它们必全为零。又由于 \(\chi_5\) 与主特征标正交，我们得到
\[
0 = (\chi_1, \chi_5) = \frac{1}{8}\left(2+(\pm2)+0+0+0\right),
\]
故 \(\chi_5(-1) = -2\). \(Q_8\) 的完整特征标表如下：
\[
\begin{array}{c|ccccc}
\text{类} & 1 & -1 & i & j & k\\ \hline
\text{大小} & 1 & 1 & 2 & 2 & 2\\
\chi_1 & 1 & 1 & 1 & 1 & 1\\
\chi_2 & 1 & 1 & 1 & -1 & -1\\
\chi_3 & 1 & 1 & -1 & 1 & -1\\
\chi_4 & 1 & 1 & -1 & -1 & 1\\
\chi_5 & 2 & -2 & 0 & 0 & 0
\end{array}
\]
（\(Q_8\) 的特征标表。）

注意 \(D_8\) 与 \(Q_8\) 有相同的特征标表，故不同构的群可能有相同的特征标表。

注意 \(Q_8\) 的 2 次表示的值也可以通过把第二正交关系应用于特征标表的每一列而容易地算出。我们把这个验证留作习题。还要注意，虽然 \(D_8\) 与 \(Q_8\) 的 2 次不可约特征标有相同的（实数值）值，\(D_8\) 的 2 次表示可由实矩阵实现，而可以证明 \(Q_8\) 在 \(\mathbb{R}\) 上没有忠实的 2 维表示（参看第 18.1 节习题 10）。

下一个例子构造 \(S_4\) 的特征标表。\(S_4\) 的共轭类由 \(1\)、\((1\ 2)\)、\((1\ 2\ 3)\)、\((1\ 2\ 3\ 4)\) 与 \((1\ 2)(3\ 4)\) 代表，大小分别为 1, 6, 8, 6 与 3. 由于 \(S_4' = A_4\)，有两个 1 次特征标：主特征标与取值等于置换符号的特征标。

为得到 2 次不可约特征标，设 \(V\) 是由 \((1\ 2)(3\ 4)\) 与 \((1\ 3)(2\ 4)\) 生成的 4 阶正规子群。\(S_4/V \cong S_3\) 的任何表示 \(\varphi\) 通过与自然投射 \(S_4 \to S_4/V\) 复合给出 \(S_4\) 的表示；若前者不可约，则后者也不可约。设 \(\varphi\) 是投射与 \(S_3\) 的不可约 2 维表示的复合，\(\chi_3\) 是它的特征标。\(S_3\) 的类 \(1\) 与 \((1\ 2)(3\ 4)\) 映到商群中的单位元，\((1\ 2)\) 与 \((1\ 2\ 3\ 4)\) 映到对换，\((1\ 2\ 3)\) 映到 3-轮换。故 \(\chi_3\) 的值可直接从 \(S_3\) 的表中 2 次特征标的值读出。

由于 \(S_4\) 有 5 个不可约特征标且次数平方和为 24，必还有两个不可约特征标，其次数都是 3. 在第 18.3 节例 2 中算出了其中一个，记作 \(\chi_4\). 回忆 \(\chi_4(\sigma)\) 等于 \(\sigma\) 的不动点个数减 1. 剩下的不可约特征标 \(\chi_5\) 是 \(\chi_4\chi_2\). 人们可以用第 18.3 节命题 17 或第 18.3 节习题 13 看出这个乘积确实是特征标，而第一正交关系验证它不可约。
\[
\begin{array}{c|ccccc}
\text{类} & 1 & (1\ 2) & (1\ 2\ 3) & (1\ 2\ 3\ 4) & (1\ 2)(3\ 4)\\ \hline
\text{大小} & 1 & 6 & 8 & 6 & 3\\
\chi_1 & 1 & 1 & 1 & 1 & 1\\
\chi_2 & 1 & -1 & 1 & -1 & 1\\
\chi_3 & 2 & 0 & -1 & 0 & 2\\
\chi_4 & 3 & 1 & 0 & -1 & -1\\
\chi_5 & 3 & -1 & 0 & 1 & -1
\end{array}
\]
（\(S_4\) 的特征标表。）

由 \(S_4\) 的特征标表可以容易地算出 \(A_4\) 的特征标表。注意 \(A_4\) 有 4 个共轭类。又 \(|A_4 : A_4'| = 3\)，故 \(A_4\) 有三个 1 次特征标，每个 1 次表示的核都含 \(V_4 = A_4'\). 剩下的不可约特征标次数必为 3. 人们直接由应用于 \(A_4\) 的正交关系验证 \(S_4\) 的特征标 \(\chi_4\) 限制到 \(A_4\)（\(= \chi_5|_{A_4}\)）不可约。这个不可约性检验确实必要，因为一个群的不可约表示限制到子群不必不可约（例如 \(S_3\) 的不可约 2 次表示限制到任何真子群时必须变得可约，因为这些真子群都阿贝尔）。\(A_4\) 的特征标表如下：
\[
\begin{array}{c|cccc}
\text{类} & 1 & (1\ 2)(3\ 4) & (1\ 2\ 3) & (1\ 3\ 2)\\ \hline
\text{大小} & 1 & 3 & 4 & 4\\
\chi_1 & 1 & 1 & 1 & 1\\
\chi_2 & 1 & 1 & \omega & \omega^2\\
\chi_3 & 1 & 1 & \omega^2 & \omega\\
\chi_4 & 3 & -1 & 0 & 0
\end{array}
\]
（\(A_4\) 的特征标表。）

其中 \(\omega\) 是 \(\mathbb{C}\) 中的本原 3 次单位根。

作为最后一个例子，我们构造 \(S_5\) 的下述特征标表：
\[
\begin{array}{c|ccccccc}
\text{类} & 1 & (1\ 2) & (1\ 2\ 3) & (1\ 2\ 3\ 4) & (1\ 2\ 3\ 4\ 5) & (1\ 2)(3\ 4) & (1\ 2)(3\ 4\ 5)\\ \hline
\text{大小} & 1 & 10 & 20 & 30 & 24 & 15 & 20\\
\chi_1 & 1 & 1 & 1 & 1 & 1 & 1 & 1\\
\chi_2 & 1 & -1 & 1 & -1 & 1 & 1 & -1\\
\chi_3 & 4 & 2 & 1 & 0 & -1 & 0 & -1\\
\chi_4 & 4 & -2 & 1 & 0 & -1 & 0 & 1\\
\chi_5 & 5 & -1 & -1 & 1 & 0 & 1 & -1\\
\chi_6 & 5 & 1 & -1 & -1 & 0 & 1 & 1\\
\chi_7 & 6 & 0 & 0 & 0 & 1 & -2 & 0
\end{array}
\]
（\(S_5\) 的特征标表。）

共轭类及其大小已在第 4.3 节算出。由于 \(|S_5 : S_5'| = 2\)，有两个 1 次特征标：主特征标与「符号」特征标。\(S_5\) 在 5 个点上的自然置换给出 5 次置换特征标。与 \(S_4\) 和 \(S_3\) 一样，正交关系表明其范数的平方为 2，且它含主特征标。故 \(\chi_3\) 是置换特征标减主特征标（并且与小对称群一样，\(\chi_3(\sigma)\) 是 \(\sigma\) 的不动点个数减 1）。如对 \(S_4\) 的论证，推出 \(\chi_4 = \chi_3\chi_2\) 也是不可约特征标。

为得到 \(\chi_5\)，回忆 \(S_5\) 有六个 Sylow 5-子群。它通过共轭作用在这些子群上给出一个忠实的 6 次置换表示。若 \(\psi\) 是相应线性表示的特征标，则由于 \(\sigma \in S_5\) 固定一个 Sylow 5-子群当且仅当它正规化该子群，有
\[
\psi(\sigma) = \sigma \text{ 所正规化的 Sylow 5-子群个数}.
\]
Sylow 5-子群 \(\langle(1\ 2\ 3\ 4\ 5)\rangle\) 在 \(S_5\) 中的正规化子是 \(\langle(1\ 2\ 3\ 4\ 5), (2\ 3\ 5\ 4)\rangle\)，而 Sylow 5-子群的所有正规化子在 \(S_5\) 中都共轭于这个群。这个正规化子只含单位元、5-轮换、4-轮换与两个不交对换之积。没有其他轮换型正规化任何 Sylow 5-子群，故在任何其他类上 \(\psi\) 为零。为计算 \(\psi\) 在剩下三个非单位类上的值，注意（在 \(S_6\) 中观察）在任何 6 个点上的忠实作用中下列事实成立：5 阶元素必是 5-轮换（故固定 1 个点）；固定一个点的 4 阶元素必是 4-轮换（故固定 2 个点）；4 阶元素的平方这个 2 阶元素也恰固定 2 个点。这给出 \(\psi\) 的全部值。现在直接计算表明
\[
\|\psi\|^2 = 2 \quad \text{且} \quad (\chi_1, \psi) = 1.
\]
故 \(\chi_5 = \psi - \chi_1\) 是不可约的 5 次特征标。由与 \(\chi_4\) 相同的理论得到 \(\chi_6 = \chi_5\chi_2\) 是另一个不可约特征标。

由于有 7 个共轭类，还剩一个不可约特征标，其次数为 6. 它的值可以立刻由正则特征标的分解得到（参看第 18.2 节例 3 与第 18.3 节例 4）：
\[
\chi_7 = \frac{\rho - \chi_1 - \chi_2 - 4\chi_3 - 4\chi_4 - 5\chi_5 - 5\chi_6}{6}.
\]
用正交关系直接计算可以验证 \(\chi_7\) 不可约。注意 \(\chi_7\) 的值是在没有显式给出具有该特征标的表示的情况下算出的。

\subsection*{习题}

\begin{enumerate}
\item 计算 \(\mathbb{Z}_2 \times \mathbb{Z}_2\)、\(\mathbb{Z}_2 \times \mathbb{Z}_3\) 与 \(\mathbb{Z}_2 \times \mathbb{Z}_2 \times \mathbb{Z}_2\) 的特征标表。解释为什么 \(\mathbb{Z}_2 \times \mathbb{Z}_3\) 的表含本原 6 次单位根。

\item 计算 \(D_{16}\) 的不可约特征标的次数。

\item 计算 \(A_5\) 的不可约特征标的次数。推出 \(S_5\) 的 6 次不可约特征标限制到 \(A_5\) 时不可约。〔\(A_5\) 的共轭类在第 4.3 节中算出。〕

\item 利用本节的各特征标表，对（a）至（d）每一种情形，用第一正交关系把指定的置换特征标写成不可约特征标之和：

（a）\(S_3\) 的子群 \(A_3\) 的置换特征标；

（b）\(S_4\) 的子群 \(\langle(1\ 2\ 3\ 4)\rangle\) 的置换特征标；

（c）\(S_4\) 的子群 \(V_4\) 的置换特征标；

（d）\(S_5\) 的子群 \(\langle(1\ 2\ 3), (1\ 2), (4\ 5)\rangle\) 的置换特征标（这个子群是 \(S_5\) 的 Sylow 3-子群的正规化子）。

\item 假设对群的任何特征标 \(\psi\)，\(\psi^2\) 也是特征标（其中 \(\psi^2(g) = (\psi(g))^2\)）——这是第 18.3 节命题 17 的特殊情形。利用本节的各特征标表，对（a）至（e）每一种情形写出指定特征标 \(\chi\) 的平方 \(\chi^2\) 的值，并用第一正交关系把 \(\chi^2\) 写成不可约特征标之和：

（a）\(\chi = \chi_3\)，\(S_3\) 表中的 2 次特征标；

（b）\(\chi = \chi_5\)，\(Q_8\) 表中的 2 次特征标；

（c）\(\chi = \chi_5\)，\(S_4\) 表中的最后一个特征标；

（d）\(\chi = \chi_4\)，\(S_5\) 表中的第二个 4 次特征标；

（e）\(\chi = \chi_7\)，\(S_5\) 表中的最后一个特征标。

\item 计算 \(A_5\) 的特征标表。

\item 证明 \(S_6\) 有一个 5 次不可约特征标。

\item 计算 \(D_{10}\) 的特征标表。（这个表含非实数的项。）

\item 计算 \(D_{12}\) 的特征标表。

\item 计算 \(S_3 \times S_3\) 的特征标表。

\item 计算 \(\mathbb{Z}_3 \times S_3\) 的特征标表。

\item 计算 \(\mathbb{Z}_2 \times S_4\) 的特征标表。

\item 计算 \(S_3 \times S_4\) 的特征标表。

\item 设 \(n \geq 3\) 是整数。证明 \(D_{2n}\) 的每个不可约特征标的次数为 1 或 2，并求出每种次数的不可约特征标个数。〔\(D_{2n}\) 的共轭类见第 4.3 节习题 31 与 32，其换位子群已在第 5.4 节算出。〕

\item 证明特征标表是可逆矩阵。〔用正交关系。〕

\item 对 \(A_5\) 与 \(D_{10}\) 分别描述哪些不可约特征标彼此代数共轭（参看第 18.3 节的习题）。

\item 设 \(p\) 是任意素数，\(P\) 是 \(p^3\) 阶非阿贝尔群（在同构意义下 \(P\) 有两种选择；对奇素数 \(p\)，它们在第 5.5 节分类 \(p^3\) 阶群时被构造出来）。本习题确定 \(P\) 的特征标表，并证明两种同构型有相同的特征标表（论证包含本节避开的 \(p = 2\) 情形）。

（a）证明 \(P\) 有 \(p^2\) 个 1 次特征标。

（b）证明 \(P\) 有 \(p - 1\) 个 \(p\) 次不可约特征标，且这些特征标连同 \(p^2\) 个 1 次特征标是 \(P\) 的全部不可约特征标。〔用第 18.2 节定理 10（3）与定理 12。〕

（c）推出（无论同构型如何）群 \(P\) 有 \(p^2+p-1\) 个共轭类，其中 \(p^2\) 个大小为 1（即是中心类），另有 \(p-1\) 个大小各为 \(p\). 还推出大小各为 \(p\) 的类恰是 \(P\) 的中心的非单位陪集（即若 \(x \in P - Z(P)\)，则 \(x\) 的共轭类是陪集 \(xZ(P)\) 中的 \(p\) 个元素之集）。

（d）证明若 \(\chi\) 是 \(p\) 次不可约特征标，则提供 \(\chi\) 的表示忠实。

（e）固定 \(P\) 的中心的一个生成元 \(z\)，并设 \(\varepsilon\) 是 \(\mathbb{C}\) 中固定的本原 \(p\) 次单位根。证明若 \(\chi\) 是 \(p\) 次不可约特征标，则 \(\chi(z) = p\varepsilon^i\)，其中 \(i \in \{1, 2, \ldots, p-1\}\). 进一步证明对 \(P - Z(P)\) 中所有 \(x\) 有 \(\chi(x) = 0\).（注意于是所有 \(p\) 次特征标彼此代数共轭。）〔用与构造 \(Q_8\) 的特征标表相同的推理。〕

（f）证明对每个 \(i \in \{1, 2, \ldots, p-1\}\) 存在唯一的 \(p\) 次不可约特征标 \(\chi_i\) 使 \(\chi_i(z) = p\varepsilon^i\). 推出 \(P\) 的特征标表被唯一确定，并给出它的描述。〔回忆第 6.1 节，无论同构型如何都有 \(P' = Z(P)\) 且 \(P/P' \cong \mathbb{Z}_p \times \mathbb{Z}_p\). 由此可以写出 1 次特征标。而（e）描述了 \(p\) 次特征标。〕
\end{enumerate}

\section{Burnside 定理与 Hall 定理}

本节给出特征标理论的一个「理论性」应用：Burnside 的 \(p^aq^b\) 定理。我们还证明 Philip Hall 对有限可解群的刻画，它的群论证明以 Burnside 定理作为归纳的第一步。

\noindent\textbf{Burnside 定理}

下述结果由 Burnside 在 1904 年证明。虽然它的纯群论证明最近才被发现（见 B. Huppert 与 N. Blackburn 的 \textit{Finite Groups III}, Springer-Verlag, 1982 中的定理 2.8），这里给出的 Burnside 原始证明非常易读、优雅且相当简短（鉴于我们现在对表示论的了解）。

\begin{theorem}\label{thm:19.1}
（\textbf{Burnside}）对素数 \(p\) 与 \(q\)，每个 \(p^aq^b\) 阶群都可解。
\end{theorem}

在着手证明 Burnside 定理本身之前，我们先建立一些一般性的结果。这些预备命题的一个容易的推论是：任何有限群的不可约特征标的次数都整除它的阶。而直接导致 Burnside 定理证明的具体结果出现在引理 6 与引理 7 中。

由此很容易推出，Burnside 定理的最小阶反例是非阿贝尔单群，而正是这两个特征标论引理通过证明一个正规子群的存在性给出矛盾。

我们首先回忆第 15.3 节代数整数的定义。

\begin{definition}
若元素 \(a \in \mathbb{C}\) 是某个系数取自 \(\mathbb{Z}\) 的首一多项式的根，则称它为\textbf{代数整数}。
\end{definition}

证明 Burnside 定理所需的基本结果是：

\begin{proposition}\label{pro:19.2}
设 \(a \in \mathbb{C}\).

（1）下列条件等价：

（i）\(a\) 是代数整数；

（ii）\(a\) 在 \(\mathbb{Q}\) 上代数，且 \(a\) 在 \(\mathbb{Q}\) 上的极小多项式有整数系数；

（iii）\(\mathbb{Z}[a]\) 是有限生成 \(\mathbb{Z}\)-模（其中 \(\mathbb{Z}[a]\) 是由 \(\mathbb{Z}\) 与 \(a\) 生成的 \(\mathbb{C}\) 的子环，即 \(a\) 的非负幂的所有 \(\mathbb{Z}\)-线性组合之环）。

（2）\(\mathbb{C}\) 中的代数整数构成一个环，而 \(\mathbb{Q}\) 中的代数整数是 \(\mathbb{Z}\) 的元素。
\end{proposition}

\begin{proof}
这些都在第 15.3 节中建立。（第 15.3 节中关于整扩张与代数整数性质的部分可以独立于第 15 章其余内容阅读。）
\end{proof}

\begin{corollary}\label{cor:19.3}
对有限群 \(G\) 的每个特征标 \(\psi\)，\(\psi(x)\) 对所有 \(x \in G\) 都是代数整数。
\end{corollary}

\begin{proof}
由第 18.3 节命题 14，\(\psi(x)\) 是单位根之和。每个单位根都是代数整数，故结果立刻由命题 2（2）得到。
\end{proof}

在开始主要证明之前，我们还需要一些预备的特征标论引理。对任意有限群 \(G\) 采用下述记号：\(\chi_1, \ldots, \chi_r\) 是 \(G\) 的互不相同的不可约（复）特征标，\(K_1, \ldots, K_r\) 是 \(G\) 的共轭类，\(\varphi_i\) 是特征标为 \(\chi_i\) 的不可约矩阵表示（对每个 \(i\)）。

\begin{proposition}\label{pro:19.4}
对每个 \(i\)，在集合 \(\{K_1, \ldots, K_r\}\) 上定义复值函数 \(\omega_i\) 为
\[
\omega_i(K_j) = \frac{|K_j|\chi_i(g)}{\chi_i(1)},
\]
其中 \(g\) 是 \(K_j\) 的任意元素。则对所有 \(i\) 与 \(j\)，\(\omega_i(K_j)\) 是代数整数。
\end{proposition}

\begin{proof}
我们先证明若 \(I\) 是单位矩阵，则
\[
\sum_{g \in K_j}\varphi_i(g) = \omega_i(K_j)I. \tag{19.1}
\]
为此设 \(X\) 是（1）的左端。如第 18.2 节所见，每个 \(x \in G\) 通过共轭作用置换 \(K_j\) 的元素，故 \(X\) 与所有 \(\varphi_i(g)\) 交换。由 Schur 引理（第 18.1 节习题 18），\(X\) 是纯量矩阵：\(X = aI\)，其中 \(a \in \mathbb{C}\). 剩下要证明 \(a = \omega_i(K_j)\). 但
\[
a\chi_i(1) = \operatorname{tr}X = |K_j|\chi_i(g),
\]
这正是建立（1）所需的。

现在固定 \(g \in K_s\)，并定义 \(a_{ijs}\) 为满足 \(g_i \in K_i\)、\(g_j \in K_j\) 且 \(g_ig_j = g\) 的有序对 \((g_i, g_j)\) 的个数。注意 \(a_{ijs}\) 是整数。它不依赖于 \(K_s\) 中 \(g\) 的选取，因为若 \(x^{-1}gx\) 是 \(g\) 的共轭，则乘积为 \(g\) 的每个有序对 \((g_i, g_j)\) 给出乘积为 \(x^{-1}gx\) 的有序对 \((x^{-1}g_ix, x^{-1}g_jx)\)（反之亦然）。

下面我们证明对所有 \(i, j, t \in \{1, \ldots, r\}\)
\[
\omega_t(K_i)\omega_t(K_j) = \sum_{s=1}^{r}a_{ijs}\omega_t(K_s). \tag{19.2}
\]
为此注意由（1），（2）的左端是下述方程左端的纯量矩阵的对角元素：
\[
\begin{aligned}
\left(\sum_{g \in K_i}\varphi_t(g)\right)\left(\sum_{g \in K_j}\varphi_t(g)\right) &= \sum_{s=1}^{r}\ \sum_{g \in K_s}a_{ijs}\varphi_t(g)\\
&= \sum_{s=1}^{r}a_{ijs}\sum_{g \in K_s}\varphi_t(g) \qquad (\text{因为 } a_{ijs} \text{ 不依赖于 } g \in K_s)\\
&= \sum_{s=1}^{r}a_{ijs}\omega_t(K_s)I \qquad (\text{由 }(1)).
\end{aligned}
\]
比较这些纯量矩阵的元素给出（2）。

现在（2）蕴含：对每个 \(t \in \{1, \ldots, r\}\)，由 \(\mathbb{Z}\) 与 \(\omega_t(K_1), \ldots, \omega_t(K_r)\) 生成的 \(\mathbb{C}\) 的子环是有限生成 \(\mathbb{Z}\)-模（它作为 \(\mathbb{Z}\)-模由 \(1, \omega_t(K_1), \ldots, \omega_t(K_r)\) 生成）。由于 \(\mathbb{Z}\) 是主理想整环，子模 \(\mathbb{Z}[\omega_t(K_j)]\) 也是有限生成 \(\mathbb{Z}\)-模，故由命题 2，\(\omega_t(K_j)\) 是代数整数。这完成了证明。
\end{proof}

\begin{corollary}\label{cor:19.5}
有限群 \(G\) 的每个复不可约表示的次数整除 \(G\) 的阶，即对 \(i = 1, 2, \ldots, r\) 有 \(\chi_i(1) \mid |G|\).
\end{corollary}

\begin{proof}
在命题 4 的记号下，取 \(g_j \in K_j\)，有
\[
\frac{|G|}{\chi_i(1)} = (\chi_i, \chi_i)\frac{|G|}{\chi_i(1)} = \sum_{j=1}^{r}\frac{|K_j|\chi_i(g_j)}{\chi_i(1)}\overline{\chi_i(g_j)} = \sum_{j=1}^{r}\omega_i(K_j)\overline{\chi_i(g_j)}.
\]
右端是代数整数，而左端是有理数，故是整数。这证明推论。
\end{proof}

下面两个引理直接导致 Burnside 定理。

\begin{lemma}\label{lem:19.6}
若群 \(G\) 有一个共轭类 \(K\) 与一个不可约矩阵表示 \(\varphi\)（特征标为 \(\chi\)）使 \((|K|, \chi(1)) = 1\)，则对 \(g \in K\)，或者 \(\chi(g) = 0\)，或者 \(\varphi(g)\) 是纯量矩阵。
\end{lemma}

\begin{proof}
由假设存在 \(s, t \in \mathbb{Z}\) 使 \(s|K|+t\chi(1) = 1\). 故
\[
s|K|\chi(g)+t\chi(1)\chi(g) = \chi(g).
\]
把两端除以 \(\chi(1)\)，并注意由推论 3 与命题 4，\(\frac{|K|\chi(g)}{\chi(1)}\) 与 \(t\chi(g)\) 都是代数整数，故 \(\frac{\chi(g)}{\chi(1)}\) 也是代数整数。设 \(a_1 = \frac{\chi(g)}{\chi(1)}\)，并设 \(a_1, a_2, \ldots, a_n\) 是它在 \(\mathbb{Q}\) 上的所有代数共轭（即 \(a_1\) 在 \(\mathbb{Q}\) 上的极小多项式的根）。由于 \(a_1\) 是单位根之和除以整数 \(\chi(1)\)，每个 \(a_i\) 也是 \(\chi(1)\) 个单位根之和除以 \(\chi(1)\). 故 \(a_i\) 的复绝对值对所有 \(i\) 都 \(\leq 1\). 现在 \(b = \prod_{i=1}^{n}a_i \in \mathbb{Q}\)，且 \(b\) 是代数整数（\(\pm b\) 是 \(a_1\) 的不可约多项式的常数项），故 \(b \in \mathbb{Z}\). 但
\[
|b| = \prod_{i=1}^{n}|a_i| \leq 1,
\]
故 \(b = 0\) 或 \(b = \pm1\). 于是或者 \(\chi(g) = 0\)，或者 \(|a_1| = 1\). 在前一种情形引理已成立，故设 \(\chi(g) \neq 0\)，从而 \(\left|\frac{\chi(g)}{\chi(1)}\right| = 1\)，即 \(|\chi(g)| = \chi(1) = n\).

设 \(\varphi_1\) 是与 \(\varphi\) 等价的矩阵表示，其中 \(\varphi_1(g)\) 是对角矩阵：
\[
\varphi_1(g) = \begin{pmatrix}\varepsilon_1 & & \\ & \ddots & \\ & & \varepsilon_n\end{pmatrix}.
\]
故 \(\chi(g) = \varepsilon_1+\cdots+\varepsilon_n\). 由三角不等式，若存在 \(i, j\) 使 \(\varepsilon_i \neq \varepsilon_j\)，则 \(|\varepsilon_1+\cdots+\varepsilon_n| < n = \chi(1)\). 由于这不可能，必有 \(\varphi_1(g) = \varepsilon I\)（其中 \(\varepsilon = \varepsilon_i\) 对所有 \(i\)）。由于纯量矩阵只与自身相似，\(\varphi(g) = \varepsilon I\) 也成立。这完成了证明。
\end{proof}

\begin{lemma}\label{lem:19.7}
若对 \(G\) 的某个非单位共轭类 \(K\)，\(|K|\) 是某个素数的幂，则 \(G\) 不是非阿贝尔单群。
\end{lemma}

\begin{proof}
反设 \(G\) 是非阿贝尔单群，并设 \(|K| = p^c\). 取 \(g \in K\). 若 \(c = 0\)，则 \(g \in Z(G)\)，与非阿贝尔单群有平凡中心矛盾。如上，设 \(\chi_1, \ldots, \chi_r\) 是 \(G\) 的全部不可约特征标，其中 \(\chi_1\) 是主特征标，\(\rho\) 是 \(G\) 的正则特征标。把 \(\rho\) 分解为不可约特征标得到
\[
0 = \rho(g) = 1+\sum_{i=2}^{r}\chi_i(1)\chi_i(g). \tag{19.3}
\]
若对所有满足 \(\chi_j(g) \neq 0\) 的 \(j > 1\) 都有 \(p \mid \chi_j(1)\)，则写 \(\chi_j(1) = pd_j\). 此时（3）成为
\[
0 = 1+p\sum_j d_j\chi_j(g).
\]
故 \(\sum_j d_j\chi_j(g) = -1/p\) 是代数整数，矛盾。这证明存在某个 \(j\) 使 \(p\) 不整除 \(\chi_j(1)\) 且 \(\chi_j(g) \neq 0\). 若 \(\varphi\) 是特征标为 \(\chi_j\) 的表示，则 \(\varphi\) 忠实（因为假设 \(G\) 单），且由引理 6，\(\varphi(g)\) 是纯量矩阵。由于 \(\varphi(g)\) 与所有矩阵交换，\(\varphi(g) \in Z(\varphi(G))\). 这迫使 \(g \in Z(G)\)，与 \(G\) 是非阿贝尔单群矛盾。引理证完。
\end{proof}

现在证明 Burnside 定理。设 \(G\) 是 \(p^aq^b\) 阶群，其中 \(p\) 与 \(q\) 是素数。若 \(p = q\)，或者某个指数为 0，则 \(G\) 幂零从而可解。故可假设不是这种情形。用归纳法，设 \(G\) 是极小阶反例。若 \(G\) 有真非平凡正规子群 \(N\)，则由归纳 \(N\) 与 \(G/N\) 都可解，从而 \(G\) 可解（参看第 3.4 节或命题 6.10）。故可假设 \(G\) 是非阿贝尔单群。设 \(P \in Syl_p(G)\). 由第 4 章定理 8，存在 \(g \in Z(P)\) 且 \(g \neq 1\). 由于 \(P \subseteq C_G(g)\)，\(g\) 的共轭类的阶（等于 \(|G : C_G(g)|\)）与 \(p\) 互素，即是 \(q\) 的幂。这与引理 7 矛盾，从而完成 Burnside 定理的证明。

\noindent\textbf{Philip Hall 定理}

回忆若有限群的子群的阶与指数互素，则称它为 \textbf{Hall 子群}。对群 \(G\) 的任何子群 \(H\)，满足 \(G = HK\) 且 \(H \cap K = 1\) 的子群 \(K\) 称为 \(H\) 在 \(G\) 中的\textbf{补}。

\begin{theorem}\label{thm:19.8}
（P. Hall）设 \(G\) 是 \(p_1^{\alpha_1}p_2^{\alpha_2}\cdots p_t^{\alpha_t}\) 阶群，其中 \(p_1, \ldots, p_t\) 是互不相同的素数。若对每个 \(i \in \{1, \ldots, t\}\) 存在 \(G\) 的子群 \(H_i\) 使 \(|G : H_i| = p_i^{\alpha_i}\)，则 \(G\) 可解。
\end{theorem}

Hall 定理也可以表述为：若对每个 \(i \in \{1, \ldots, t\}\)，\(G\) 的一个 Sylow \(p_i\)-子群有补，则 \(G\) 可解。Hall 定理的逆也成立——这就是第 6.1 节习题 33.

我们首先需要一些初等引理。

\begin{lemma}\label{lem:19.9}
若 \(G\) 可解且 \(|G| > 1\)，则存在 \(P \trianglelefteq G\)，其中 \(P\) 是某个素数 \(p\) 的非平凡 \(p\)-群。
\end{lemma}

\begin{proof}
这是第 6.1 节末关于可解群的极小正规子群的习题的特殊情形。容易看出这一点：设 \(P\) 是 \(G\) 的换位子群列中最后一个非平凡项 \(G^{(n-1)}\) 的非平凡 Sylow 子群（其中 \(G\) 的可解长度为 \(n\)）。此时 \(G^{(n-1)}\) 阿贝尔，故 \(P\) 是 \(G^{(n-1)}\) 的特征子群，从而是 \(G\) 的正规子群。
\end{proof}

\begin{lemma}\label{lem:19.10}
设 \(G\) 是 \(p_1^{\alpha_1}p_2^{\alpha_2}\cdots p_t^{\alpha_t}\) 阶群，其中 \(p_1, \ldots, p_t\) 是互不相同的素数。假设存在 \(G\) 的子群 \(H\) 与 \(K\)，使对每个 \(i \in \{1, \ldots, t\}\)，或者 \(p_i^{\alpha_i}\) 整除 \(|H|\)，或者 \(p_i^{\alpha_i}\) 整除 \(|K|\). 则 \(G = HK\)，且 \(|H \cap K| = (|H|, |K|)\).
\end{lemma}

\begin{proof}
固定某个 \(i \in \{1, \ldots, t\}\)，先设 \(p_i^{\alpha_i}\) 整除 \(H\) 的阶。由于 \(HK\) 是 \(H\) 的右陪集的不交并，而每个右陪集的阶等于 \(|H|\)，推出 \(p_i^{\alpha_i}\) 整除 \(|HK|\). 类似地，若 \(p_i^{\alpha_i}\) 整除 \(|K|\)，由于 \(HK\) 是 \(K\) 的左陪集的不交并，同样有 \(p_i^{\alpha_i}\) 整除 \(|HK|\). 故 \(|G| \mid |HK|\)，从而 \(G = HK\). 由于
\[
|HK| = \frac{|H||K|}{|H \cap K|},
\]
推出 \(|H \cap K| = (|H|, |K|)\).
\end{proof}

我们现在开始证明 Hall 定理，对 \(|G|\) 归纳。注意若 \(t = 1\)，则假设对任何群都平凡地成立（\(H_1 = 1\)）；而若 \(t = 2\)，则对任何群假设也成立（由 Sylow 定理，\(H_1\) 是 \(G\) 的 Sylow \(p_2\)-子群，而 \(H_2\) 是 \(G\) 的 Sylow \(p_1\)-子群）。若 \(t = 1\)，则 \(G\) 幂零，从而可解；若 \(t = 2\)，则由 Burnside 定理 \(G\) 可解。故假设 \(t \geq 3\).

固定 \(i\)，注意由上一个引理，对所有 \(j \in \{1, \ldots, t\} - \{i\}\)，
\[
|H_i \cap H_j| = \frac{|H_i|}{p_j^{\alpha_j}}.
\]
故 \(H_i\) 的每个 Sylow \(p_j\)-子群在 \(H_i\) 中都有补：\(H_i \cap H_j\). 由归纳，\(H_i\) 可解。

由引理 9，我们可以选取 \(P \trianglelefteq H_1\)，使 \(|P| = p_i^{\alpha_i} > 1\) 对某个 \(i > 1\) 成立。由于 \(t \geq 3\)，存在指标 \(j \in \{1, \ldots, t\} - \{1, i\}\). 由引理 10
\[
|H_1 \cap H_j| = p_2^{\alpha_2}\cdots p_{j-1}^{\alpha_{j-1}}p_{j+1}^{\alpha_{j+1}}\cdots p_t^{\alpha_t}.
\]
故 \(H_1 \cap H_j\) 含 \(H_1\) 的一个 Sylow \(p_i\)-子群。由于 \(P\) 是 \(H_1\) 的正规 \(p_i\)-子群，\(P\) 含于 \(H_1\) 的每个 Sylow \(p_i\)-子群，故 \(P \subseteq H_1 \cap H_j\). 由引理 10，\(G = H_1H_j\)，故每个 \(g \in G\) 可写成 \(g = h_1h_j\)，其中 \(h_1 \in H_1\)、\(h_j \in H_j\). 于是
\[
gPg^{-1} = h_1h_jPh_j^{-1}h_1^{-1} \subseteq h_1H_jh_1^{-1},
\]
由此 \(G\) 中 \(H_j\) 的所有共轭之交等于 \(\bigcap_{h_1 \in H_1}h_1H_jh_1^{-1}\). 现在由于 \(P \subseteq H_j\) 且 \(P \trianglelefteq H_1\)，对每个 \(h_1 \in H_1\) 有
\[
P = h_1Ph_1^{-1} \subseteq h_1H_jh_1^{-1}.
\]
故
\[
1 \neq P \subseteq \bigcap_{h_1 \in H_1}h_1H_jh_1^{-1} = \bigcap_{g \in G}gH_jg^{-1} = N.
\]
于是 \(N\) 是 \(G\) 的非平凡正规子群；又 \(N \subseteq H_j \neq G\)，故 \(N\) 是真子群。由此 \(N\) 与 \(G/N\) 都满足定理的假设（参看第 3.3 节的习题）。由归纳 \(N\) 与 \(G/N\) 都可解，故 \(G\) 可解。这完成了 Hall 定理的证明。

\subsection*{习题}

\begin{enumerate}
\item 证明对称群 \(S_n\) 的每个特征标都是整数值的（即对所有 \(g \in S_n\) 与 \(S_n\) 的所有特征标 \(\psi\) 有 \(\psi(g) \in \mathbb{Z}\)）。〔见第 18.3 节习题 22。〕

\item 设 \(G\) 是有限群，且具有性质：每个极大子群的指数或者是素数，或者是素数的平方。证明 \(G\) 可解。（单群 \(GL_3(\mathbb{F}_2)\) 具有性质：每个极大子群的指数或者是 7，或者是 8，即或者是素数，或者是素数的立方——参看第 6.2 节。）〔设 \(p\) 是整除 \(|G|\) 的最大素数，\(P\) 是 \(G\) 的 Sylow \(p\)-子群。若 \(P \trianglelefteq G\)，则对 \(G/P\) 用归纳。否则设 \(M\) 是包含 \(N_G(P)\) 的极大子群。用第 4.5 节习题 51 证明 \(p = 3\)，并推出 \(|G| = 2^a3^b\).  〕

\item 假设有限群 \(G\) 有一个指数是某个素数的幂的阿贝尔子群 \(H\)。证明 \(G\) 可解。

\item 把上一习题中的「阿贝尔」换成「幂零」，重复该习题。

\item 用 Philip Hall 定理证明中的想法，在不使用特征标理论的情况下证明 Burnside 的 \(p^aq^b\) 定理在所有 Sylow 子群都阿贝尔的特殊情形。
\end{enumerate}

\section{诱导特征标理论引论}

设 \(G\) 是有限群，\(H\) 是 \(G\) 的子群，\(\varphi\) 是子群 \(H\) 在任意域 \(F\) 上的表示。本节我们说明如何从子群的表示得到一个 \(G\) 的表示，称为诱导表示。我们还确定这个诱导表示的特征标（诱导特征标）用 \(\varphi\) 的特征标表示的公式，并通过计算具体群中的一些诱导特征标来说明这个公式。最后，我们应用诱导特征标理论证明不存在 \(3^3\cdot 7\cdot 13\cdot 409\) 阶单群——这个群阶在第 6.2 节末讨论单群存在性问题时出现过。诱导表示与诱导特征标理论标志着更高等的表示论的开端。本节意在作为引论而非全面的论述，我们收录的结果都是为此目的而选取的。

首先注意，可能无法把一个子群 \(H\) 的表示延拓为 \(G\) 的表示 \(\psi\) 使 \(\psi|_H = \varphi\). 例如设 \(A_3 \leq S_3\)，而 \(A_3\) 有一个忠实的 1 次表示（参看第 1 节）。由于 \(S_3\) 的每个 1 次表示的核都含 \(A_3 = S_3'\)，这个 \(A_3\) 的表示不能延拓为 \(S_3\) 的表示。另一个不能延拓到整个群的子群表示的例子如下：取 \(G\) 为任意单群，\(\varphi\) 为 \(H\) 的任意表示且具有性质「\(\ker\varphi\) 是 \(H\) 的真非平凡正规子群」。若 \(\varphi\) 延拓为 \(G\) 的表示 \(\psi\)，则 \(\psi\) 的核将是 \(G\) 的真非平凡子群，与 \(G\) 单矛盾。我们将看到，诱导特征标的方法从子群 \(H\) 的给定表示 \(\varphi\) 产生 \(G\) 的表示 \(\psi\)，但一般地 \(\psi|_H \neq \varphi\)（事实上，除非 \(H = G\)，\(\psi\) 的次数将大于 \(\varphi\) 的次数）。

我们在第 10.4 节推论 9 后的例 5 中看到，由于 \(FH\) 是 \(FG\) 的子环，环 \(FG\) 是 \((FG, FH)\)-双模；故对任何左 \(FH\)-模 \(V\)，阿贝尔群 \(FG \otimes_{FH}V\) 是左 \(FG\)-模（称为 \(V\) 从 \(FH\) 到 \(FG\) 的标量扩张）。在有限群的表示论中，这个扩张有一个专门的名称。

\begin{definition}
设 \(H\) 是有限群 \(G\) 的子群，\(V\) 是提供 \(H\) 的表示 \(\varphi\) 的 \(FH\)-模。则 \(FG\)-模 \(FG \otimes_{FH}V\) 称为 \(V\) 的\textbf{诱导模}，它所提供的 \(G\) 的表示称为 \(\varphi\) 的\textbf{诱导表示}。若 \(\psi\) 是 \(\varphi\) 的特征标，则诱导表示的特征标称为\textbf{诱导特征标}，记作 \(\operatorname{Ind}_H^G(\psi)\).
\end{definition}

\begin{theorem}\label{thm:19.11}
设 \(H\) 是有限群 \(G\) 的子群，\(g_1, \ldots, g_m\) 是 \(G\) 中 \(H\) 的互不相同的左陪集的代表元。设 \(V\) 是提供 \(H\) 的 \(n\) 次矩阵表示 \(\varphi\) 的 \(FH\)-模。则 \(FG\)-模 \(W = FG \otimes_{FH}V\) 在 \(F\) 上的维数为 \(nm\)，且 \(W\) 有一组基，使 \(W\) 提供由
\[
\Psi(g) = \begin{pmatrix}
\varphi(g_1^{-1}gg_1) & \varphi(g_1^{-1}gg_2) & \cdots & \varphi(g_1^{-1}gg_m)\\
\varphi(g_2^{-1}gg_1) & \varphi(g_2^{-1}gg_2) & \cdots & \varphi(g_2^{-1}gg_m)\\
\vdots & \vdots & & \vdots\\
\varphi(g_m^{-1}gg_1) & \varphi(g_m^{-1}gg_2) & \cdots & \varphi(g_m^{-1}gg_m)
\end{pmatrix}
\]
定义的对每个 \(g \in G\) 的矩阵表示 \(\Psi\)，其中每个 \(\varphi(g_i^{-1}gg_j)\) 是出现在 \(\Psi(g)\) 的第 \(i,j\) 分块位置的 \(n \times n\) 分块，而当 \(g_i^{-1}gg_j \notin H\) 时定义 \(\varphi(g_i^{-1}gg_j)\) 为零分块。
\end{theorem}

\begin{proof}
首先注意 \(FG\) 是自由右 \(FH\)-模：
\[
FG = g_1FH \oplus g_2FH \oplus \cdots \oplus g_mFH.
\]
由于张量积与直和交换（第 10.4 节定理 17），作为阿贝尔群有
\[
W = FG \otimes_{FH}V \cong (g_1 \otimes V) \oplus (g_2 \otimes V) \oplus \cdots \oplus (g_m \otimes V).
\]
由于 \(F\) 在 \(FG\) 的中心中，推出这也是 \(F\)-向量空间同构。故若 \(v_1, v_2, \ldots, v_n\) 是提供矩阵表示 \(\varphi\) 的 \(V\) 的一组基，则 \(\{g_i \otimes v_j \mid 1 \leq i \leq m,\ 1 \leq j \leq n\}\) 是 \(W\) 的一组基。这表明 \(W\) 的维数为 \(mn\). 把这组基按 \(m\) 组、每组 \(n\) 个排序为由 \(\{g_1 \otimes v_1, \ldots, g_1 \otimes v_n\}, \ldots, \{g_m \otimes v_1, \ldots, g_m \otimes v_n\}\) 给出的 \(m\) 个块。

我们计算每个 \(g\) 关于这组基作用在 \(W\) 上的矩阵表示 \(\Psi(g)\). 固定 \(j\) 与 \(g\)，设 \(gg_j = g_ih\)，其中 \(i\) 是某个指标、\(h \in H\). 则对每个 \(k\)
\[
g(g_j \otimes v_k) = g_jh \otimes v_k = g_j \otimes hv_k = g_j \otimes \left(\sum_{l=1}^{n}a_{lk}(h)v_l\right) = \sum_{l=1}^{n}a_{lk}(h)(g_j \otimes v_l),
\]
其中 \(a_{lk}(h)\) 是 \(h\) 关于基 \(\{v_1, \ldots, v_n\}\) 作用在 \(V\) 上的矩阵的 \(l,k\) 系数。换句话说，\(g\) 在 \(W\) 上的作用把 \(W\) 的第 \(j\) 块 \(n\) 个基向量映到第 \(i\) 块基向量，然后在该块上具有矩阵 \(\varphi(h)\). 由于 \(h = g_i^{-1}gg_j\)，这就描述了定理中的分块矩阵 \(\Psi(g)\)，正如所需。
\end{proof}

\begin{corollary}\label{cor:19.12}
按定理 11 的记号，

（1）若 \(\psi\) 是 \(V\) 提供的特征标，则诱导特征标由
\[
\operatorname{Ind}_H^G(\psi)(g) = \sum_{i=1}^{m}\psi(g_i^{-1}gg_i)
\]
给出，其中当 \(g_i^{-1}gg_i \notin H\) 时定义 \(\psi(g_i^{-1}gg_i) = 0\)；以及

（2）若 \(g\) 在 \(G\) 中不共轭于 \(H\) 的任何元素，则 \(\operatorname{Ind}_H^G(\psi)(g) = 0\). 特别地，若 \(H\) 是 \(G\) 的正规子群，则 \(\operatorname{Ind}_H^G(\psi)\) 在所有 \(G - H\) 的元素上为零。
\end{corollary}

\noindent\textbf{注：}由于 \(H\) 的特征标 \(\psi\) 在 \(H\) 的共轭类上是常值，对所有 \(h \in H\) 有 \(\psi(g) = \psi(h^{-1}gh)\). 当 \(h\) 取遍 \(H\) 的所有元素时，\(xh\) 取遍陪集 \(xH\) 的所有元素。故诱导特征标的公式也可以写成
\[
\operatorname{Ind}_H^G(\psi)(g) = \frac{1}{|H|}\sum_{x \in G}\psi(x^{-1}gx),
\]
其中每个固定陪集中的各元素给出相同的特征标值 \(|H|\) 次（这正是因子 \(1/|H|\) 的来源），并且同样当 \(x^{-1}gx \notin H\) 时 \(\psi(x^{-1}gx) = 0\).

\begin{proof}
由上面算出的 \(g\) 的矩阵，\(\Psi(g)\) 对角线上各分块 \(\varphi(g_i^{-1}gg_i)\) 除 \(g_i^{-1}gg_i \in H\) 外为零。故分块矩阵 \(\Psi(g)\) 的迹是满足 \(g_i^{-1}gg_i \in H\) 的矩阵 \(\varphi(g_i^{-1}gg_i)\) 的迹之和。由于 \(\varphi(g_i^{-1}gg_i)\) 的迹是 \(\psi(g_i^{-1}gg_i)\)，（1）成立。

若对所有陪集代表元 \(g_i\) 都有 \(g_i^{-1}gg_i \notin H\)，则 \(\operatorname{Ind}_H^G(\psi)(g)\) 的和中每项为零。特别地，若 \(g\) 不在正规子群 \(H\) 中，则 \(g\) 的任何共轭也不在 \(H\) 中，故 \(\operatorname{Ind}_H^G(\psi)\) 在 \(g\) 上为零。
\end{proof}

\noindent\textbf{例}

（1）设 \(G = D_{12} = \langle r, s \mid r^6 = s^2 = 1,\ rs = sr^{-1}\rangle\) 是 12 阶二面体群，\(H = \{1, r^3, s, sr^3\}\)，故 \(H\) 同构于 Klein 四元群且 \(|G : H| = 3\). 按定理 11 的记号，我们给出一个具体的 \(H\) 的表示下诱导表示在 \(r\) 与 \(s\) 上的矩阵。设 \(H\) 在 \(\mathbb{Q}\) 上 2 维向量空间关于某组基 \(v_1, v_2\) 的表示由
\[
\varphi(s) = \begin{pmatrix}-1 & 0\\ 0 & 1\end{pmatrix} = A, \quad \varphi(r^3) = \begin{pmatrix}1 & 0\\ 0 & 1\end{pmatrix} = B, \quad \varphi(sr^3) = \begin{pmatrix}-1 & 0\\ 0 & -1\end{pmatrix} = C
\]
给出，故 \(n = 2\)、\(m = 3\)，诱导表示 \(\Phi\) 的次数为 \(nm = 6\). 固定 \(H\) 在 \(G\) 中左陪集的代表元 \(g_1 = 1\)、\(g_2 = r\) 与 \(g_3 = r^2\)，故 \(g_k = r^{k-1}\). 则
\[
g_i^{-1}rg_j = r^{-(i-1)}rr^{j-1} = r^{j-i+1}, \qquad \text{而} \qquad g_i^{-1}sg_j = sr^{i+j-2}.
\]
故诱导表示的 \(6 \times 6\) 矩阵为
\[
\Phi(r) = \begin{pmatrix}0 & 0 & B\\ I & 0 & 0\\ 0 & I & 0\end{pmatrix}, \qquad \Phi(s) = \begin{pmatrix}A & 0 & 0\\ 0 & 0 & C\\ 0 & C & 0\end{pmatrix},
\]
其中 \(2 \times 2\) 矩阵 \(A\)、\(B\)、\(C\) 如上给出，\(I\) 是 \(2 \times 2\) 单位矩阵，而 \(0\) 表示 \(2 \times 2\) 零矩阵。

（2）若 \(H\) 是 \(G\) 的任意子群，\(\psi_1\) 是 \(H\) 的主特征标，则 \(\operatorname{Ind}_H^G(\psi_1)(g)\) 对每个满足 \(g_i^{-1}gg_i \in H\) 的陪集代表元 \(g_i\) 计 1. 由于 \(g_i^{-1}gg_i \in H\) 当且仅当 \(g\) 在左乘作用下固定左陪集 \(g_iH\)，\(\operatorname{Ind}_H^G(\psi_1)(g)\) 就是 \(g\) 在 \(H\) 的左陪集上的置换表示中不动点的个数。故由第 18.3 节例 3 我们看到：若 \(\psi_1\) 是 \(H\) 的主特征标，则 \(\operatorname{Ind}_H^G(\psi_1)\) 是 \(H\) 在 \(G\) 中的左陪集上的置换特征标。在 \(H = 1\) 的特殊情形，这蕴含：若 \(\chi_1\) 是平凡子群 \(H = 1\) 的主特征标，则 \(\operatorname{Ind}_1^G(\chi_1)\) 是 \(G\) 的正则特征标。这也表明一般地，即使被诱导的特征标不可约，诱导特征标也不必不可约。

（3）设 \(G = S_3\)，\(\psi\) 是 \(A_3 = \langle x\rangle\) 的非主 1 次特征标，故对某个本原 3 次单位根 \(\zeta\) 有 \(\psi(x) = \zeta\)（\(A_3 = \mathbb{Z}_3\) 与 \(S_3\) 的特征标表见第 1 节）。设 \(\Psi = \operatorname{Ind}_{A_3}^{S_3}(\psi)\). 于是 \(\Psi\) 的次数为 \(1 \cdot |S_3 : A_3| = 2\)，且由推论，\(\Psi\) 在所有对换上为零。若 \(y\) 是任意对换，则 \(1, y\) 是 \(A_3\) 在 \(S_3\) 中的一组左陪集代表元，且 \(y^{-1}xy = x^2\). 故 \(\Psi(x) = \psi(x)+\psi(x^2) = \zeta+\zeta^2 = -1\). 这表明若 \(\psi\) 是 \(A_3\) 的两个非主不可约特征标中的任意一个，则 \(\psi\) 的诱导特征标是 \(S_3\) 的（唯一）2 次不可约特征标。特别地，子群的不同特征标可能诱导出整个群的同一个特征标。

（4）设 \(G = D_8\) 有通常的生成元与关系，\(H = \langle s\rangle\). 设 \(\psi\) 是 \(H\) 的非主不可约特征标，\(\Psi = \operatorname{Ind}_H^G(\psi)\). 取 \(H\) 的左陪集代表元 \(1, r, r^2, r^3\). 由定理 11，\(\Psi(1) = 4\). 由于 \(\psi(s) = -1\)，直接计算得 \(\Psi(s) = -2\). 由推论 12（2）得到 \(\Psi(r) = \Psi(r^2) = \Psi(sr) = 0\). 在第 1 节 \(D_8\) 特征标表的记号下，由正交关系得到 \(\Psi = \chi_2+\chi_4+\chi_5\)（这可以通过观察检验）。

本节余下部分取域 \(F\) 为复数：\(F = \mathbb{C}\).

在用诱导特征标应用于单群之前，我们先计算一类重要群的特征标。

\begin{definition}
若有限群 \(G\) 有一个真非平凡正规子群 \(Q\)，且对 \(Q\) 的所有非单位元 \(x\) 有 \(C_G(x) \subseteq Q\)，则称 \(G\) 是以 \(Q\) 为 \textbf{Frobenius 核}的 \textbf{Frobenius 群}。
\end{definition}

鉴于本节开头提到的对单群的应用，我们把注意力限制在 \(q^ap\) 阶的 Frobenius 群 \(G\)，其中 \(p\) 与 \(q\) 是不同的素数，Frobenius 核 \(Q\) 是 \(q^a\) 阶初等阿贝尔 \(q\)-群，且循环群 \(G/Q\) 通过共轭不可约地作用在 \(Q\) 上。换句话说，我们假设 \(Q\) 是 \(q\) 阶循环群的直积，且含于 \(Q\) 中的 \(G\) 的正规子群只有 1 与 \(Q\)，即 \(Q\) 是 \(G\) 的极小正规子群。例如，\(A_4\) 是这种类型的 Frobenius 群，其 Frobenius 核是 \(V_4\)（它的 Sylow 2-子群）。又如，若 \(p\) 与 \(q\) 是不同素数且 \(p < q\)，\(G\) 是 \(pq\) 阶非阿贝尔群（当 \(p \mid q-1\) 时总存在），则 \(G\) 是以其 Sylow \(q\)-子群为 Frobenius 核的 Frobenius 群（由 Sylow 定理它是正规的）。我们将基本上确定这些 Frobenius 群的特征标表。关于更一般的 Frobenius 群的类似结果出现在习题中。

\begin{proposition}\label{pro:19.13}
设 \(G\) 是 \(q^ap\) 阶 Frobenius 群，其中 \(p\) 与 \(q\) 是不同素数，Frobenius 核 \(Q\) 是 \(q^a\) 阶初等阿贝尔 \(q\)-群，且循环群 \(G/Q\) 通过共轭不可约地作用在 \(Q\) 上。则下列结论成立：

（1）\(G = QP\)，其中 \(P\) 是 \(G\) 的 Sylow \(p\)-子群。\(G\) 的每个非单位元的阶为 \(p\) 或 \(q\). 每个 \(p\) 阶元素都共轭于 \(P\) 的某个元素，而每个 \(q\) 阶元素都属于 \(Q\). \(P\) 的非单位元代表 \(p-1\) 个互不相同的 \(p\) 阶元素共轭类，每个这样的类大小为 \(q^a\). 有 \((q^a-1)/p\) 个互不相同的 \(q\) 阶元素共轭类，每个这样的类大小为 \(p\).

（2）\(G' = Q\)，故 \(G\) 的 1 次特征标个数为 \(p\)，且每个 1 次特征标的核都含 \(Q\).

（3）若 \(\psi\) 是 \(Q\) 的任意非主不可约特征标，则 \(\operatorname{Ind}_Q^G(\psi)\) 是 \(G\) 的不可约特征标。此外，\(G\) 的每个次数 \(> 1\) 的不可约特征标都等于某个 \(Q\) 的非主不可约特征标 \(\psi\) 的 \(\operatorname{Ind}_Q^G(\psi)\). \(G\) 的每个不可约特征标的次数为 1 或 \(p\)，而次数为 \(p\) 的不可约特征标个数为 \((q^a-1)/p\).
\end{proposition}

\begin{proof}
注意由阶的考虑有 \(QP = G\). 按 Frobenius 群的定义，又由于 \(Q\) 阿贝尔，对 \(Q\) 的每个非单位元 \(h\) 有 \(C_G(h) = Q\). 若 \(x\) 是 \(pq\) 阶元素，则 \(x^p\) 是 \(q\) 阶元素，故属于 \(G\) 的唯一个 Sylow \(q\)-子群 \(Q\). 但此时 \(x\) 与 \(x^p\) 交换，故 \(x \in C_G(x^p) = Q\)，矛盾。故 \(G\) 没有 \(pq\) 阶元素。

由 Sylow 定理，每个 \(p\) 阶元素都共轭于 \(P\) 的某个元素，每个 \(q\) 阶元素都属于 \(Q\). \(P\) 的两个不同元素在 \(G\) 中不共轭，因为若对某个 \(x, y \in P\) 有 \(g^{-1}xg = y\)，则在阿贝尔群 \(G/Q\) 中有 \(\bar x = \bar y\)，即 \(x \equiv y \pmod Q\). 又 \(x, y \in P\) 且 \(P \cap Q = 1\)，故 \(x = y\). 于是恰有 \(p-1\) 个 \(p\) 阶元素的共轭类，它们由 \(P\) 的非单位元代表。若 \(x\) 是 \(P\) 的非单位元，则 \(C_G(x) = P\)，故 \(x\) 的共轭类含 \(|G : P| = q^a\) 个元素。最后，若 \(h\) 是 \(Q\) 的非单位元，则 \(C_G(h) = Q\)，而 \(h\) 的共轭类是 \(\{h, h^x, \ldots, h^{x^{p-1}}\}\)，其中 \(P = \langle x\rangle\). 这证明（1）的所有部分。

由于 \(G/Q\) 阿贝尔有 \(G' \subseteq Q\). 又 \(G\) 非阿贝尔，且按假设 \(Q\) 是 \(G\) 的极小正规子群，必有 \(G' = Q\). （2）现在由第 18.2 节推论 11 得到。

设 \(\psi\) 是 \(Q\) 的非主不可约特征标，并令 \(\Psi = \operatorname{Ind}_Q^G(\psi)\). 我们用正交关系证明 \(\Psi\) 不可约。设 \(1, x, \ldots, x^{p-1}\) 是 \(Q\) 在 \(G\) 中的陪集代表元。由推论 12（2），\(\Psi\) 在 \(G - Q\) 上为零，故
\[
\|\Psi\|^2 = \frac{1}{|G|}\sum_{h \in Q}\Psi(h)\overline{\Psi(h)}
= \frac{1}{|G|}\sum_{h \in Q}\sum_{i,j=0}^{p-1}\psi(x^ihx^{-i})\overline{\psi(x^jhx^{-j})}
= \frac{p}{|G|}\sum_{h \in Q}\psi(h)\overline{\psi(h)}
= \frac{p|Q|}{|G|} = 1,
\]
其中第二行由诱导特征标 \(\Psi\) 的定义得到，第三行因为 \(Q\) 的每个元素在第二行的和中恰好出现 \(p\) 次（这一步用到 \(G/Q\) 在 \(Q\) 上的作用不可约，故非主特征标 \(\psi\) 在 \(P\) 的作用下的稳定子为 1），最后一行由 \(Q\) 中的第一正交关系得到，因为 \(\psi\) 是 \(Q\) 的不可约特征标。这证明 \(\Psi\) 是 \(G\) 的不可约特征标。

下面证明 \(G\) 的每个次数 \(> 1\) 的不可约特征标都是 \(Q\) 的某个非主 1 次特征标的诱导特征标，方法是计算这样得到的 \(G\) 的互不相同的不可约特征标个数。由（1）与（2），\(G\) 的不可约特征标个数（即共轭类数）为 \(p+(q^a-1)/p\)，而 1 次特征标个数为 \(p\). 故 \(G\) 的次数 \(> 1\) 的不可约特征标个数为 \((q^a-1)/p\). 群 \(P\) 如下作用在 \(Q\) 的非主不可约特征标之集 \(C\) 上：对每个 \(\psi \in C\) 与每个 \(x \in P\)，定义 \(\psi^x\) 为
\[
\psi^x(h) = \psi(xhx^{-1}) \qquad \text{对所有 } h \in Q.
\]
由于 \(\psi\) 是从 \(Q\) 到 \(\mathbb{C}^\times\) 的非平凡同态（回忆阿贝尔群 \(Q\) 的所有不可约特征标都是 1 次的），容易推出 \(\psi^x\) 也是同态。故 \(\psi^x \in C\)，于是 \(P\) 置换 \(C\) 的元素。现在设 \(x\) 是循环群 \(P\) 的生成元，则 \(1, x, \ldots, x^{p-1}\) 是 \(Q\) 在 \(G\) 中的左陪集代表元。由推论 12 应用于这组陪集代表元，我们看到：若 \(\psi \in C\)，则 \(\operatorname{Ind}_Q^G(\psi)\) 在 \(Q\) 的任何元素 \(h\) 上的值由和 \(\psi(h)+\psi^x(h)+\cdots+\psi^{x^{p-1}}(h)\) 给出。因此当诱导特征标 \(\operatorname{Ind}_Q^G(\psi)\) 限制到 \(Q\) 时，它分解为 \(Q\) 的不可约特征标之直和：
\[
\operatorname{Ind}_Q^G(\psi)|_Q = \psi+\psi^x+\cdots+\psi^{x^{p-1}}.
\]
若 \(\psi_1\) 与 \(\psi_2\) 处于 \(P\) 作用在 \(C\) 上的不同轨道，则诱导特征标 \(\operatorname{Ind}_Q^G(\psi_1)\) 与 \(\operatorname{Ind}_Q^G(\psi_2)\) 限制到 \(Q\) 是互不相同的特征标（它们没有公共的不可约成分）。故由 \(P\) 在 \(C\) 的不同轨道中的元素诱导得到的特征标是 \(G\) 的互不相同的次数为 \(p\) 的不可约特征标。阿贝尔群 \(Q\) 有 \(q^a-1\) 个非主不可约特征标（即 \(|C| = q^a-1\)），而 \(|P| = p\)，故 \(P\) 在 \(C\) 上至少有 \((q^a-1)/p\) 个轨道，从而至少有这么多 \(G\) 的次数为 \(p\) 的互不相同的不可约特征标。由于 \(G\) 恰有 \((q^a-1)/p\) 个次数 \(> 1\) 的不可约特征标，\(G\) 的每个次数 \(> 1\) 的不可约特征标其次数必为 \(p\)，且必是某个 \(C\) 中元素的诱导特征标。证明完成。
\end{proof}

对最后一个例子，我们需要诱导特征标的两条性质。这两条性质列在下一个命题中，其证明是直接的习题，容易由诱导特征标的公式或由诱导模的定义连同张量积的性质得到。

\begin{proposition}\label{pro:19.14}
设 \(G\) 是群，\(H\) 是 \(G\) 的子群，\(\psi\) 与 \(\psi'\) 是 \(H\) 的特征标。

（1）（特征标的诱导是加性的）\(\operatorname{Ind}_H^G(\psi+\psi') = \operatorname{Ind}_H^G(\psi)+\operatorname{Ind}_H^G(\psi')\).

（2）（特征标的诱导是传递的）若 \(H \leq K \leq G\)，则
\[
\operatorname{Ind}_K^G(\operatorname{Ind}_H^K(\psi)) = \operatorname{Ind}_H^G(\psi).
\]
\end{proposition}

由命题 14 的（1）推出：若 \(\sum_{i=1}^{t}n_i\psi_i\) 是 \(H\) 的特征标的任意整线性组合，其中对所有 \(i\) 有 \(n_i \geq 0\)，则
\[
\operatorname{Ind}_H^G\left(\sum_{i=1}^{t}n_i\psi_i\right) = \sum_{i=1}^{t}n_i\operatorname{Ind}_H^G(\psi_i). \tag{*}
\]
\(H\) 的形如 \(\sum_{i=1}^{t}n_i\psi_i\) 的类函数（其中系数是任意整数，不必非负）称为 \(H\) 的\textbf{广义特征标}或\textbf{虚特征标}。对 \(H\) 的广义特征标，我们用等式（*）定义它在 \(G\) 中的诱导广义特征标（此时允许负系数）。这样，\(\operatorname{Ind}_H^G\) 成为从 \(H\) 的广义特征标的加法群到 \(G\) 的广义特征标的加法群的群同态（它把特征标映到特征标）。这意味着推论 12 中诱导特征标的公式在 \(\psi\) 是 \(H\) 的广义特征标时也成立。

\noindent\textbf{对 \(3^3\cdot 7\cdot 13\cdot 409\) 阶群的应用}

下面我们以证明下述结果作为本节的结束：

\noindent 不存在 \(3^3\cdot 7\cdot 13\cdot 409\) 阶单群。

正如本节开头提到的，这种阶的单群在第 6.2 节末讨论单群存在性问题时出现过。可以用涉及一个 819 次置换表示的论证证明不存在这种阶的单群（参看第 6.2 节的习题）。我们在此给出一个特征标论的证明，因为这些方法说明了有限群论中一些重要的想法。这个途径基于 M. Suzuki 的开创性论文 \textit{The nonexistence of a certain type of simple group of odd order}，Proc. Amer. Math. Soc., 8(1957), pp. 686--695，该文处理远为一般的群。由于我们处理的是具体的群阶，我们的论证更简单、数值上更明确，但仍保留 Suzuki 工作中的一些关键想法。此外，Suzuki 的论文及其后继者 W. Feit、M. Hall 与 J. Thompson 的 \textit{Finite groups in which the centralizer of any non-identity element is nilpotent}，Math. Zeit., 74(1960), pp. 1--17，是冗长而困难的 Feit--Thompson 定理（参看第 3.4 节）的原型。我们的讨论也传达了这些奠基性论文的某些风格。特别地，这些论文中的每一篇都遵循这样的基本发展：先确定 Sylow 子群的结构与嵌入方式，然后应用特征标理论（大量依赖诱导特征标）。

本节余下部分假设 \(G\) 是 \(3^3\cdot 7\cdot 13\cdot 409\) 阶单群。我们列出 \(G\) 的一些可以用第 6.2 节中源于 Sylow 定理的方法验证的性质。细节留作习题。

（1）设 \(q_1 = 3\)，\(Q_1\) 是 \(G\) 的 Sylow 3-子群，\(N_1 = N_G(Q_1)\). 则 \(Q_1\) 是 \(3^3\) 阶初等阿贝尔 3-群，\(N_1\) 是 \(3^3\cdot 13\) 阶 Frobenius 群，其 Frobenius 核为 \(Q_1\)，且 \(N_1/Q_1\) 通过共轭不可约地作用在 \(Q_1\) 上。

（2）设 \(q_2 = 7\)，\(Q_2\) 是 \(G\) 的 Sylow 7-子群，\(N_2 = N_G(Q_2)\). 则 \(Q_2\) 是 7 阶循环群，\(N_2\) 是 \(7\cdot 3\) 阶非阿贝尔群（故 \(N_2\) 是以 \(Q_2\) 为 Frobenius 核的 Frobenius 群）。

（3）设 \(q_3 = 13\)，\(Q_3\) 是 \(G\) 的 Sylow 13-子群，\(N_3 = N_G(Q_3)\). 则 \(Q_3\) 是 13 阶循环群，\(N_3\) 是 \(13\cdot 3\) 阶非阿贝尔群（故 \(N_3\) 是以 \(Q_3\) 为 Frobenius 核的 Frobenius 群）。

（4）设 \(q_4 = 409\)，\(Q_4\) 是 \(G\) 的 Sylow 409-子群，\(N_4 = N_G(Q_4)\). 则 \(Q_4\) 是 409 阶循环群，\(N_4\) 是 \(409\cdot 3\) 阶非阿贝尔群（故 \(N_4\) 是以 \(Q_4\) 为 Frobenius 核的 Frobenius 群）。

（5）\(G\) 的每个非单位元的阶都是素数，且对每个 \(i = 1, 2, 3, 4\) 与每个 \(g \in G - N_i\) 有 \(Q_i \cap Q_i^g = 1\). \(G\) 的非单位共轭类为：

（a）2 个 3 阶元素的类（每个类大小为 \(7\cdot 13\cdot 409\)）；

（b）2 个 7 阶元素的类（每个类大小为 \(3^3\cdot 13\cdot 409\)）；

（c）4 个 13 阶元素的类（每个类大小为 \(3^3\cdot 7\cdot 409\)）；

（d）136 个 409 阶元素的类（每个类大小为 \(3^3\cdot 7\cdot 13\)），

故 \(G\) 有 145 个共轭类。

由于每个群 \(N_i\) 都是满足命题 13 假设的 Frobenius 群，\(N_i\) 的次数 \(> 1\) 的特征标个数可以从该命题读出：

（i）\(N_1\) 有 2 个 13 次不可约特征标；

（ii）\(N_2\) 有 2 个 3 次不可约特征标；

（iii）\(N_3\) 有 4 个 3 次不可约特征标；

（iv）\(N_4\) 有 136 个 3 次不可约特征标。

从现在起，为简化记号，对 \(G\) 的任何子群 \(H\) 与 \(H\) 的任何广义特征标 \(\mu\)，令
\[
\mu^* = \operatorname{Ind}_H^G(\mu),
\]
故星号总表示从子群到整个群 \(G\) 的诱导，而子群可以从上下文看出。

下述引理是证明的关键。它说明推论 12 中描述的诱导特征标的消失性（连同 Sylow 子群 \(Q_i\) 的平凡交性质，即对所有 \(g \in G - N_G(Q_i)\) 有 \(Q_i \cap Q_i^g = 1\)）如何用来把某些广义特征标的内积与它们的诱导广义特征标的内积联系起来。在这些计算中，广义特征标在单位元上取零这一点很重要（这也解释了为什么我们考虑同次数的特征标之差）。

\begin{lemma}\label{lem:19.15}
对任意 \(i \in \{1,2,3,4\}\)，设 \(q = q_i\)、\(Q = Q_i\)、\(N = N_i\)，并设 \(p = |N : Q|\). 设 \(\psi_1, \psi_2, \psi_3, \psi_4\) 是 \(N\) 的任意四个次数为 \(p\) 的不可约特征标（不必互不相同），并设 \(\alpha = \psi_1-\psi_2\) 与 \(\beta = \psi_3-\psi_4\). 则 \(\alpha\) 与 \(\beta\) 是 \(N\) 的广义特征标，它们在 \(N\) 的每个阶不等于 \(q\) 的元素上为零。此外，\(\alpha^*\) 与 \(\beta^*\) 是 \(G\) 的广义特征标，它们在 \(G\) 的每个阶不等于 \(q\) 的元素上为零，并且
\[
(\alpha^*, \beta^*)_G = (\alpha, \beta)_N
\]
（其中 \((\ ,\ )_H\) 表示群 \(H\) 中计算的类函数的通常 Hermite 积）。换句话说，从 \(N\) 到 \(G\) 的诱导在这类 \(N\) 的广义特征标 \(\alpha, \beta\) 上是保内积的映射。
\end{lemma}

\begin{proof}
由命题 13，存在 \(Q\) 的非主 1 次特征标 \(\lambda_1, \ldots, \lambda_4\) 使 \(\psi_j = \operatorname{Ind}_Q^N(\lambda_j)\)（\(j = 1, \ldots, 4\)）。故由推论 12，每个 \(\psi_j\) 在 \(N - Q\) 上为零，从而 \(\alpha\) 与 \(\beta\) 也如此。注意由于对所有 \(j\) 有 \(\psi_j(1) = p\)，我们有 \(\alpha(1) = \beta(1) = 0\). 由诱导的传递性，对所有 \(j\) 有 \(\psi_j^* = \operatorname{Ind}_Q^G(\lambda_j)\). 再次由推论 12 应用于后一个诱导特征标，我们看到 \(\psi_j^*\) 在所有不共轭于 \(Q\) 的某个元素的 \(G\) 的元素上为零，从而 \(\alpha^*\) 与 \(\beta^*\) 也都如此。由于诱导特征标 \(\psi_j^*\) 都有次数 \(|G : Q|\)，广义特征标 \(\alpha^*\) 与 \(\beta^*\) 在单位元上为零。故 \(\alpha^*\) 与 \(\beta^*\) 在 \(G\) 的所有阶不等于 \(q\) 的元素上为零。

最后，若 \(g_1, \ldots, g_m\) 是 \(N\) 在 \(G\) 中的左陪集代表元且 \(g_1 = 1\)，则由于对所有 \(k > 1\) 有 \(Q \cap Q^{g_k} = 1\)（由上面的（5）），由诱导（广义）特征标的公式立刻推出：对 \(Q\) 的所有非单位元 \(x\)（即 \(N\) 的所有 \(q\) 阶元素 \(x\)）有 \(\alpha^*(x) = \alpha(x)\) 与 \(\beta^*(x) = \beta(x)\). 进一步，由 Sylow 定理，\(G\) 的每个 \(q\) 阶元素都在 \(Q\) 的某个共轭中，故集合 \(Q-\{1\}\) 的 \(G\)-共轭之集把 \(G\) 中的 \(q\) 阶元素划分成 \(|G:N|\) 个不交子集。由于 \(\alpha^*\) 与 \(\beta^*\) 是 \(G\) 上的类函数，\(\alpha^*(x)\) 与 \(\beta^*(x)\) 之积在任何一个这样的子集上的和相同。这些事实蕴含
\[
(\alpha^*, \beta^*)_G = \frac{1}{|G|}\sum_{x \in G}\alpha^*(x)\overline{\beta^*(x)}
= \frac{1}{|G|}\sum_{\substack{x \in G\\ |x| = q}}\alpha^*(x)\overline{\beta^*(x)}
= \frac{1}{|G|}|G:N|\sum_{\substack{x \in N\\ |x| = q}}\alpha(x)\overline{\beta(x)}
= \frac{1}{|N|}\sum_{x \in N}\alpha(x)\overline{\beta(x)} = (\alpha, \beta)_N.
\]
这完成了证明。
\end{proof}

下一个引理在 \(N_i\) 的次数 \(> 1\) 的不可约特征标与 \(G\) 的某些非主不可约特征标之间建立对应。

\begin{lemma}\label{lem:19.16}
对任意 \(i \in \{1,2,3,4\}\)，设 \(q = q_i\)、\(Q = Q_i\)、\(N = N_i\)，并设 \(p = |N : Q|\). 设 \(\psi_1, \ldots, \psi_k\) 是 \(N\) 的互不相同的 \(p\) 次不可约特征标。则存在 \(G\) 的互不相同的不可约特征标 \(\chi_1, \ldots, \chi_k\)（它们都有相同的次数）与一个固定的符号 \(\varepsilon = \pm1\)，使对所有 \(j = 2, 3, \ldots, k\) 有
\[
\psi_1^*-\psi_j^* = \varepsilon(\chi_1-\chi_j).
\]
\end{lemma}

\begin{proof}
设 \(\alpha_j = \psi_1-\psi_j\)（\(j = 2, 3, \ldots, k\)），故 \(\alpha_j\) 满足引理 15 的假设。由于 \(\psi_1 \neq \psi_j\)，由引理 15
\[
2 = \|\alpha_j\|_N^2 = (\alpha_j, \alpha_j)_N = (\alpha_j^*, \alpha_j^*)_G = \|\alpha_j^*\|_G^2
\]
对所有 \(j\) 成立。故 \(\alpha_j^*\) 必须以 \(G\) 的两个互不相同的不可约特征标为不可约成分。由于 \(\alpha_j(1) = 0\)，它必是两个互不相同的不可约特征标之差，而这两个特征标有相同的次数。特别地，若 \(k = 2\)（这正是 \(q = 3\) 与 \(q = 7\) 的情形），引理成立。

故假设 \(k > 2\)，并写
\[
\alpha_2^* = \psi_1^*-\psi_2^* = \varepsilon(\chi-\chi'), \qquad \alpha_3^* = \psi_1^*-\psi_3^* = \varepsilon_1(\theta-\theta'),
\]
其中 \(\chi, \chi', \theta, \theta'\) 是 \(G\) 的不可约特征标，\(\varepsilon, \varepsilon_1 \in \{\pm1\}\). 如上所证，\(\chi \neq \chi'\) 且 \(\theta \neq \theta'\). 必要时交换 \(\theta\) 与 \(\theta'\)，可以假定 \(\varepsilon = \varepsilon_1\). 于是
\[
\alpha_2^*-\alpha_3^* = \psi_2^*-\psi_3^* = (\psi_2-\psi_3)^* = \varepsilon(\theta-\theta'-\chi+\chi').
\]
由引理 15，\(\psi_2^*-\psi_3^* = (\psi_2-\psi_3)^*\) 也恰有两个互不相同的不可约成分，故或者 \(\theta = \chi\)，或者 \(\theta' = \chi'\). 必要时把 \(\varepsilon\) 换成 \(-\varepsilon\)，可以假定 \(\theta = \chi\)，于是现在有
\[
\alpha_2^* = \psi_1^*-\psi_2^* = \varepsilon(\chi-\chi'), \qquad \alpha_3^* = \psi_1^*-\psi_3^* = \varepsilon(\chi-\theta'),
\]
其中 \(\chi\)、\(\chi'\) 与 \(\theta'\) 是 \(G\) 的互不相同的不可约特征标，而符号 \(\varepsilon\) 已确定。把 \(\chi\) 记为 \(\chi_1\)、\(\chi'\) 记为 \(\chi_2\)、\(\theta'\) 记为 \(\chi_3\). 现在同样可以验证：对每个 \(j \geq 3\) 存在 \(G\) 的不可约特征标 \(\chi_j\)，使
\[
\alpha_j^* = \psi_1^*-\psi_j^* = \varepsilon(\chi_1-\chi_j),
\]
且 \(\chi_1, \ldots, \chi_k\) 互不相同。由于所有 \(\chi_j\) 与 \(\chi_1\) 有相同次数，证明完成。
\end{proof}

我们指出，未必有 \(\chi_j = \psi_j^*\)，而只是 \(N\) 的不可约特征标之差诱导为 \(G\) 的不可约特征标之差。

由引理 16 得到的 \(G\) 的不可约特征标 \(\chi_i\) 称为与 \(Q\) 相联系的\textbf{例外特征标}。

\begin{lemma}\label{lem:19.17}
与 \(Q_i\) 相联系的例外特征标与与 \(Q_j\) 相联系的例外特征标全都不同，其中 \(i\) 与 \(j\) 是 \(\{1,2,3,4\}\) 的不同元素。
\end{lemma}

\begin{proof}
设 \(\chi\) 是与 \(Q_i\) 相联系的例外特征标，\(\theta\) 是与 \(Q_j\) 相联系的例外特征标。由构造，存在 \(N_i\) 的广义特征标 \(\alpha\) 与 \(N_j\) 的广义特征标 \(\beta\)，它们是 \(N_i\)（分别地 \(N_j\)）的两个次数为 \(|N_i:Q_i|\)（分别地 \(|N_j:Q_j|\)）的特征标之差，且使 \(\alpha^* = \chi-\chi'\) 与 \(\beta^* = \theta-\theta'\)，其中 \(\chi, \chi'\) 与 \(\theta, \theta'\) 是 \(G\) 的不可约特征标。由引理 15，\(\alpha^*\) 在 \(G\) 的所有阶不等于 \(q_i\) 的元素上为零（包括单位元），而 \(\beta^*\) 在 \(G\) 的所有阶不等于 \(q_j\) 的元素上为零。由于 \(q_i \neq q_j\)，显然
\[
(\alpha^*, \beta^*) = 0.
\]
由此容易推出 \(\alpha^*\) 的两个不可约成分与 \(\beta^*\) 的两个不可约成分两两正交。这建立引理。
\end{proof}

现在容易证明这样的单群 \(G\) 不存在。由引理 16 与 \(G\) 的性质（i）至（iv），我们可以数出例外特征标的个数：

（i）与 \(Q_1\) 相联系的例外特征标有 2 个；

（ii）与 \(Q_2\) 相联系的例外特征标有 2 个；

（iii）与 \(Q_3\) 相联系的例外特征标有 4 个；

（iv）与 \(Q_4\) 相联系的例外特征标有 136 个。

用 \(d_i\) 表示与 \(Q_i\) 相联系的例外特征标的公共次数（\(i = 1, \ldots, 4\)）。由引理 17，这些例外特征标给出 144 个非主不可约特征标，故它们连同主特征标就是 \(G\) 的全部不可约特征标（\(G\) 的共轭类数为 145）。不可约特征标次数的平方和为群的阶：
\[
1+2d_1^2+2d_2^2+4d_3^2+136d_4^2 = 1004913.
\]
化简得
\[
d_1^2+d_2^2+2d_3^2+68d_4^2 = 502456. \tag{19.4}
\]
最后，由于单群 \(G\) 的每个非主不可约表示都忠实，且 \(N_1\) 的忠实表示的最小次数为 13，所以每个 \(d_i \geq 13\). 又 \(d_4 < \sqrt{502456/68} < 86\)，且 \(d_4\) 整除 \(|G|\)，故
\[
d_4 \in \{13, 21, 27, 39, 63\}.
\]
进一步，由推论 5，每个 \(d_i\) 整除 \(|G|\)，故每个 \(d_i\) 只有很少几种可能。现在验证方程（4）没有解（用计算机做这件事尤其容易）。这个矛盾完成了证明。

\subsection*{习题}

本节的习题中所有表示都在复数上。

\begin{enumerate}
\item 设 \(G = S_3\)，\(H = A_3\)，\(V\) 是提供 \(A_3\) 的自然置换表示的 3 维 \(\mathbb{C}H\)-模。更确切地说，设 \(V\) 有基 \(e_1, e_2, e_3\)，并设 \(\sigma \in A_3\) 通过 \(\sigma e_i = e_{\sigma(i)}\) 作用在 \(V\) 上。设 \(1\) 与 \((1\ 2)\) 是 \(A_3\) 在 \(S_3\) 中的左陪集代表元，并写出定理 11 中描述的 \(S_3\) 在诱导模 \(W\) 上的作用对 \(S_3\) 的每个元素给出的显式矩阵。

\item 在（a）至（f）的每一种情形中，指定了某个群 \(G\) 的子群 \(H\) 的一个特征标 \(\psi\). 计算诱导特征标 \(\operatorname{Ind}_H^G(\psi)\) 在 \(G\) 的所有共轭类上的值，并用第 1 节的特征标表把 \(\operatorname{Ind}_H^G(\psi)\) 写成不可约特征标之和：

（a）\(\psi\) 是 \(S_3\) 的子群 \(\langle(1\ 2)\rangle\) 的唯一非主 1 次特征标；

（b）\(\psi\) 是 \(D_8\) 的子群 \(\langle r\rangle\) 的由 \(\psi(r) = i\) 定义的 1 次特征标，其中 \(i \in \mathbb{C}\) 是 \(-1\) 的平方根；

（c）\(\psi\) 是 \(D_8\) 的子群 \(\langle r\rangle\) 的由 \(\psi(r) = -1\) 定义的 1 次特征标；

（d）\(\psi\) 是 \(S_4\) 的子群 \(V_4 = \langle(1\ 2), (3\ 4)\rangle\) 的任意非主 1 次特征标；

（e）\(\psi = \chi_4\) 是第 1 节中 \(H = S_4\) 的特征标表的两个 3 次特征标中的第一个，而 \(H\) 是 \(G = S_5\) 的子群；

（f）\(\psi\) 是 \(S_5\) 的子群 \(V_4 = \langle(1\ 2), (3\ 4)\rangle\) 的任意非主 1 次特征标。

\item 用命题 13 显式写出下列每个群的特征标表：

（a）10 阶二面体群；

（b）57 阶非阿贝尔群；

（c）有一个正规初等阿贝尔 Sylow 2-子群的 56 阶非阿贝尔群。

\item 设 \(H\) 是 \(G\) 的子群，\(\varphi\) 是 \(H\) 的表示，并假设 \(N\) 是 \(G\) 的正规子群且 \(N \leq H\)、\(N\) 含于 \(\varphi\) 的核中。证明 \(N\) 也含于 \(\varphi\) 的诱导表示的核中。

\item 设 \(N\) 是 \(G\) 的正规子群，\(\psi_1\) 是 \(N\) 的主特征标。设 \(\psi\) 是诱导特征标 \(\operatorname{Ind}_N^G(\psi_1)\)，故由上一习题我们可以把 \(\psi\) 看作 \(G/N\) 的某个表示的特征标。证明 \(\psi\) 是 \(G/N\) 的正则表示的特征标。

\item 设 \(Z\) 是 \(G\) 的中心的任意子群，\(|G : Z| = m\)，\(\psi\) 是 \(Z\) 的特征标。证明
\[
\operatorname{Ind}_Z^G(\psi)(g) = \begin{cases}m\psi(g), & \text{若 } g \in Z\\ 0, & \text{若 } g \notin Z.\end{cases}
\]

\item 设 \(\varphi\) 是 \(G\) 的子群 \(H\) 的矩阵表示，并对每个 \(g \in G\) 用定理 11 陈述中的显示公式定义矩阵 \(\Phi(g)\). 通过证明对所有 \(x, y \in G\) 有 \(\Phi(xy) = \Phi(x)\Phi(y)\) 来直接证明 \(\Phi\) 是表示。

\item 设 \(G\) 是以 \(Q\) 为 Frobenius 核的 Frobenius 群。假设 \(Q\) 与 \(G/Q\) 都阿贝尔，但 \(G\) 不阿贝尔（即 \(G \neq Q\)）。设 \(|Q| = n\)，\(|G : Q| = m\).

（a）证明 \(G/Q\) 循环，并证明 \(G = QC\)，其中 \(C\) 是 \(G\) 的某个循环子群且 \(C \cap Q = 1\)（即 \(G\) 是 \(Q\) 与 \(C\) 的半直积，且 \(|C| = m\)）。〔设 \(q\) 是 \(n\) 的一个素因子，并设 \(G/Q\) 通过共轭作用在初等阿贝尔 \(q\)-群 \(\{h \in Q \mid h^q = 1\}\) 上。对一个不可约成分应用第 18.1 节习题 14（e）与 Frobenius 群的定义。〕

（b）证明 \(n\) 与 \(m\) 互素。〔若素数 \(p\) 同时整除 \(Q\) 的阶与指数，设 \(P\) 是 \(G\) 的 Sylow \(p\)-子群。则 \(P \cap Q \trianglelefteq P\) 且 \(P \cap Q\) 是 \(Q\) 的 Sylow \(p\)-子群。考虑子群 \(Z(P) \cap Q\) 在 \(G\) 中的中心化子（由第 6.1 节定理 1，这个交非平凡）。〕

（c）证明 \(G\) 没有 \(qp\) 阶元素，其中 \(q\) 是 \(n\) 的任何非平凡因子，\(p\) 是 \(m\) 的任何非平凡因子。〔如命题 13 那样论证。〕

（d）证明 \(G\) 中属于 \(Q\) 的非单位共轭类个数为 \((n-1)/m\)，且每个这样的类大小为 \(m\).〔如命题 13 那样论证。〕

（e）证明 \(C\) 的两个不同元素在 \(G\) 中不共轭。推出 \(C\) 的非单位元代表 \(m-1\) 个互不相同的 \(G\) 的共轭类，且每个这样的类大小为 \(n\). 再推出 \(G-Q\) 的每个元素都共轭于 \(C\) 的某个元素，且 \(G\) 有 \(m+(n-1)/m\) 个共轭类。

（f）证明 \(G' = Q\)，并推出 \(G\) 有 \(m\) 个互不相同的 1 次特征标。〔为证明 \(Q \leq G'\)，设 \(C = \langle x\rangle\)，并论证映射 \(h \mapsto [h,x] = x^{-1}h^{-1}xh\) 是从 \(Q\) 到 \(Q\) 的同态，其核平凡，故这个映射是满射。〕

（g）证明若 \(\psi\) 是 \(Q\) 的任意非主不可约特征标，则 \(\operatorname{Ind}_Q^G(\psi)\) 是 \(G\) 的不可约特征标。证明 \(G\) 的每个次数 \(>1\) 的不可约特征标都等于某个 \(Q\) 的非主不可约特征标 \(\psi\) 的 \(\operatorname{Ind}_Q^G(\psi)\). 推出 \(G\) 的每个不可约特征标的次数为 1 或 \(m\)，而次数为 \(m\) 的不可约特征标个数为 \((n-1)/m\).〔验证命题 13（3）的证明在适当修改所涉及的数值后建立了这个更一般的结果。〕

\item 用上一习题显式写出 \(\langle(1\ 2\ 3\ 4\ 5), (2\ 3\ 5\ 4)\rangle\) 的特征标表，它是 \(S_5\) 中 Sylow 5-子群的正规化子（这个群是 20 阶 Frobenius 群）。

\item 设 \(N\) 是 \(G\) 的正规子群，\(\psi\) 是 \(N\) 的特征标，\(g \in G\). 定义 \(\psi^g\) 为 \(\psi^g(h) = \psi(ghg^{-1})\)（对所有 \(h \in N\)）。

（a）证明 \(\psi^g\) 是 \(N\) 的特征标（\(\psi\) 与 \(\psi^g\) 称为 \(N\) 的 \textbf{\(G\)-共轭}特征标）。证明 \(\psi^g\) 不可约当且仅当 \(\psi\) 不可约。

（b）证明映射 \(\psi \mapsto \psi^g\) 是 \(G\) 在 \(N\) 的特征标之集上的右群作用，且 \(N\) 在这个作用的核中。

（c）证明若 \(\psi_1\) 与 \(\psi_2\) 是 \(N\) 的 \(G\)-共轭特征标，则 \(\operatorname{Ind}_N^G(\psi_1) = \operatorname{Ind}_N^G(\psi_2)\). 还证明若 \(\psi_1\) 与 \(\psi_2\) 是 \(N\) 的不 \(G\)-共轭的特征标，则 \(\operatorname{Ind}_N^G(\psi_1) \neq \operatorname{Ind}_N^G(\psi_2)\).〔用命题 13（3）证明中的论证。〕

\item 证明若 \(G = A_4\) 而 \(N = V_4\) 是它的 Sylow 2-子群，则 \(N\) 的任意两个非主不可约特征标都是 \(G\)-共轭的（参看上一习题）。

\item 设 \(G = D_{2n}\) 由通常的生成元与关系给出。证明若 \(\psi\) 是 \(H = \langle r\rangle\) 的任意 1 次特征标且 \(\psi \neq \psi^s\)，则 \(\operatorname{Ind}_H^G(\psi)\) 是 \(D_{2n}\) 的不可约特征标。证明 \(D_{2n}\) 的每个不可约特征标都是 \(\langle r\rangle\) 的某个 1 次特征标的诱导特征标。

\item 证明命题 14 的两个部分。

\item 证明下述结果，它称为 \textbf{Frobenius 互反律}：设 \(H < G\)，\(\psi\) 是 \(H\) 的任意特征标，\(\chi\) 是 \(G\) 的任意特征标。则
\[
(\psi, \chi|_H)_H = (\operatorname{Ind}_H^G(\psi), \chi)_G.
\]
〔用诱导特征标 \(\operatorname{Ind}_H^G(\psi)\) 的公式展开右端，或仿照第 17.2 节 Shapiro 引理的证明。〕

\item 假设 \(G\) 是 \(3^3\cdot 7\cdot 13\cdot 409\) 阶单群，其 Sylow 子群及它们的正规化子由本节的性质（1）至（5）描述。证明由 \(G\) 作用在子群 \(N_4\) 的左陪集上所得的 819 次置换特征标分解为 \(\chi_0+\chi+\chi'\)，其中 \(\chi_0\) 是 \(G\) 的主特征标，而 \(\chi\) 与 \(\chi'\) 是 \(G\) 的两个互不相同的 409 次不可约特征标。〔用第 18.3 节习题 9 证明这个置换特征标 \(\pi\) 有 \(\|\pi\|^2 = 3\).〕
\end{enumerate}
